%
% @author  Andre Tampubolon, Andreas Febrian
% @version 2.2.0
% @edit by Muhammad Aulia Adil Murtito, Ichlasul Affan
%

\chapter*{\uppercase{Halaman Pernyataan Persetujuan Publikasi\\Tugas Akhir untuk Kepentingan Akademis}}
\vspace*{-1cm}
\par\noindent\rule{\textwidth}{3pt}
\vspace*{1cm}
\noindent
Sebagai sivitas akademik Universitas Indonesia,~\ifx\blank\npmDua{saya}\else{kami}\fi~yang bertanda
tangan di bawah ini:

\vspace*{0.2cm}

\pagestyle{onlypage}

\begin{tabular}{p{4.2cm} l p{6.5cm}}
	Nama&: & \penulisSatu
		\ifx\blank\npmDua\else{~/~\penulisDua}\fi
		\ifx\blank\npmTiga\else{~/~\penulisTiga}\fi \\
	NPM&: & \npmSatu
		\ifx\blank\npmDua\else{~/~\npmDua}\fi
		\ifx\blank\npmTiga\else{~/~\npmTiga}\fi \\
	Program Studi&: & \programSatu
		\ifx\blank\npmDua\else{~/~\programDua}\fi
		\ifx\blank\npmTiga\else{~/~\programTiga}\fi \\
	Jenis Karya & : & \type \\
\end{tabular}

\vspace*{0.2cm}

\noindent demi pengembangan ilmu pengetahuan, menyetujui untuk memberikan
kepada Universitas Indonesia \bo{Hak Bebas Royalti Noneksklusif
(\textit{Non-exclusive Royalty Free Right})} atas karya ilmiah~\ifx\blank\npmDua{saya}\else{kami}\fi~yang berjudul:
\begin{center}
	\judul
\end{center}
beserta perangkat yang ada (jika diperlukan). Dengan Hak Bebas Royalti
Noneksklusif ini Universitas Indonesia berhak menyimpan,
mengalihmedia/formatkan, mengelola dalam bentuk pangkalan data
(\textit{database}), merawat, dan memublikasikan tugas akhir~\ifx\blank\npmDua{saya}\else{kami}\fi~selama
tetap mencantumkan nama~\ifx\blank\npmDua{saya}\else{kami}\fi~sebagai penulis/pencipta dan sebagai
pemilik Hak Cipta.

\noindent Demikian pernyataan ini~\ifx\blank\npmDua{saya}\else{kami}\fi~buat dengan sebenarnya.

% Untuk input gambar tanda tangan, silahkan sesuaikan xshift, yshift, dan width dengan gambar tanda tangan Anda
%\begin{tikzpicture}[remember picture,overlay,shift={(current page.north east)}]
%\node[anchor=north east,xshift=-9cm,yshift=-23.5cm]{\includegraphics[width=3cm]{assets/pics/tanda_tangan_wikipedia.png}};
%\end{tikzpicture}


\begin{center}
	\vspace*{0.8cm}
	\begin{tabular}{lll}
		Dibuat di&: & Depok \\
		Pada tanggal&: & \tanggalSiapSidang \\
	\end{tabular}\\

	\vspace*{0.2cm}
	Yang menyatakan \\
	\ifx\blank\npmDua
		\vspace*{2cm}
		(\penulisSatu) \\
	\else
		\ifx\blank\npmTiga
			\begin{multicols}{2}
				\vspace*{0.5cm}
				(\penulisSatu) \\

				\vspace*{0.5cm}
				(\penulisDua) \\
			\end{multicols}
		\else
			\begin{multicols}{3}
				\vspace*{0.5cm}
				(\penulisSatu) \\

				\vspace*{0.5cm}
				(\penulisDua) \\

				\vspace*{0.5cm}
				(\penulisTiga) \\
			\end{multicols}
		\fi
	\fi
\end{center}

\newpage
